% Transcription littérale du TD 05.
% Les abréviations et l'ordre des raisonnements de la copie sont conservés.

\exercise{9}{}

\begin{statementbox}
Soit $(u_n)_{n \in \N} \in \R_+^\N / \lim_{n \to +\infty} u_n \sum_{k=0}^n u_k^2 = 1$. Mq $u_n \sim \frac{1}{\sqrt[3]{3n}}$. A-t-on une réciproque ?
\end{statementbox}

\begin{solution}
Pour $n \in \N$, posons $S_n = \sum_{k=0}^n u_k^2$.
On a $(S_n)_n$ croissante et $u_n^2 = S_n - S_{n-1}$. Par hypothèse, $\lim_{n \to +\infty} (S_n - S_{n-1}) S_n^2 = 1$.
Comme $(S_n)_n$ est croissante elle admet une limite $l \in \R_+ \cup \{+\infty\}$.
Si $l$ est finie, $\lim_{n \to +\infty} u_n = \frac{1}{l} > 0$ d'où $(S_n)_n$ diverge grossièrement, ce qui n'est pas.
Donc $l = +\infty$.
De plus $S_n^3 - S_{n-1}^3 = (S_n - S_{n-1})(S_n^2 + S_n S_{n-1} + S_{n-1}^2)$.
Et $S_n = S_{n-1} + u_n^2 = S_{n-1} + o(S_{n-1})$ car $\lim S_{n-1} = +\infty$, ainsi $S_{n-1} \sim S_n$.
Alors $S_n^3 - S_{n-1}^3 = (S_n - S_{n-1}) \underbrace{S_n^2}_{\sim S_n^2} + (S_n - S_{n-1}) \underbrace{S_n S_{n-1}}_{\sim S_n^2} + (S_n - S_{n-1}) \underbrace{S_{n-1}^2}_{\sim S_n^2}$.
Or on a $\lim_{n \to +\infty} (S_n - S_{n-1}) S_n^2 = 1$ on a donc $\lim_{n \to +\infty} S_n^3 - S_{n-1}^3 = 3$.
Par le théorème de Cesàro $\lim_{n \to +\infty} \frac{S_n^3}{n} = 3$. D'où $S_n \sim \sqrt[3]{3n}$ et $u_n \sim \frac{1}{\sqrt[3]{3n}}$.

Réciproquement si $u_n \sim \frac{1}{\sqrt[3]{3n}}$. On pose $f : [1, +\infty[ \to \R, t \mapsto \frac{1}{3 \sqrt[3]{t^2}}$ décroissant.
Pour $n \geq 1$, soit $t \in [n, n+1]$, $f(n+1) \leq f(t) \leq f(n)$.
D'où $f(n+1) \leq \int_n^{n+1} f(t) dt \leq f(n)$.
Soit $N > 1$, on somme pour $n \in \llbracket 1, N \rrbracket$ :
D'où : $\int_1^{N+1} \frac{dt}{t^{2/3}} \leq \sum_{k=1}^N \frac{1}{k^{2/3}} \leq \int_1^N \frac{dt}{t^{2/3}} + 1$.
$\sim 3(N+1)^{1/3} \sim 3N^{1/3}$ et $[3t^{1/3}]_1^N \sim 3N^{1/3}$.
D'où $\sum_{k=1}^N \frac{1}{k^{2/3}} \sim 3N^{1/3}$.
Par sommation des relations de comparaison (si $u_n \sim \frac{1}{\sqrt[3]{3n}} > 0$ tg de SATP dv) :
$\sum_{k=0}^n u_k^2 \sim \frac{1}{\sqrt[3]{9}} \times 3 n^{1/3} = \sqrt[3]{3n}$.
D'où $\lim_{n \to +\infty} u_n \sum_{k=0}^n u_k^2 = 1$.

\begin{remarkbox}
Rem : On a $(S_n - S_{n-1}) S_n^2 \to 1$ d'où :
$S_n - S_{n-1} \sim \frac{1}{S_n^2} \to 0 : S_n \sim S_{n-1}$ et $S_n - S_{n-1} \sim \frac{1}{S_{n-1}^2} + o\left(\frac{1}{S_{n-1}^2}\right)$.
On cherche $\alpha \in \R / S_n^\alpha - S_{n-1}^\alpha$ ait un équivalent "simple".
On a :
\begin{align*}
S_n^\alpha &= S_{n-1}^\alpha \left(1 + \frac{1}{S_{n-1}^3}\right)^{\alpha/3} = S_{n-1}^\alpha \left( 1 + \alpha \frac{1}{S_{n-1}^3} + o\left(\frac{1}{S_{n-1}^3}\right) \right) \\
&= S_{n-1}^\alpha + \alpha S_{n-1}^{\alpha-3} + o(S_{n-1}^{\alpha-3})
\end{align*}
Pour $\alpha = 3 : S_n^3 - S_{n-1}^3 \to 3$.
\end{remarkbox}
\end{solution}

\exercise{25}{}

\begin{statementbox}
Lundi 5/10

Soit $u_0 \in \left]0; 1\right[$ ; on définit par récurrence $\forall n \in \N, u_{n+1} = u_n - u_n^2$
\begin{enumerate}[label=\alph*.]
  \item Étudier $(u_n)_{n \in \N}$.
  \item On pose $\forall n \in \N, v_n = \frac{1}{u_n}$ ; mq $v_n = n + \ln n + \mathcal{O}(1)$.
  \item En déduire un D.A à 2 termes de $u_n$.
  \item Que dire si $u_0 \notin \left]0; 1\right[$ ?
\end{enumerate}
\end{statementbox}

\begin{solution}
\partanswer{a.}
Soit $f : [0; 1] \longrightarrow \R, x \longmapsto x - x^2$ de classe $\mathcal{C}^1$.
$\forall x \in [0; 1], f'(x) = 1 - 2x$.
On déduit du tableau de variation que $f([0; 1]) \subset [0; 1]$.
\begin{center}
\begin{tabular}{|c|lcccr|}
\hline
$x$ & $0$ & & $\frac{1}{2}$ & & $1$ \\
\hline
$f'(x)$ & & $+$ & $0$ & $-$ & \\
\hline
$f(x)$ & $0$ & $\nearrow$ & $1/4$ & $\searrow$ & $0$ \\
\hline
\end{tabular}
\end{center}
Et donc par récurrence immédiate, $\forall n \in \N, u_n \in [0; 1]$.
Et $\forall n \in \N, u_{n+1} = u_n - u_n^2 \leq u_n$ donc $(u_n)$ décroit, et est minorée par 0 : elle \boxed{CV} vers $l \in [0; 1]$.
Par continuité de $f, l = f(l)$
$l = l - l^2$ d'où $\boxed{l = 0}$.

\partanswer{b.}
Par récurrence $\forall n \in \N, u_n \in \left]0; 1\right]$, donc $v_n = \frac{1}{u_n}$ est bien défini.
On a : $\forall n \in \N, v_{n+1} - v_n = \frac{1}{u_{n+1}} - \frac{1}{u_n} = \frac{1}{u_n - u_n^2} - \frac{1}{u_n}$
Or $\lim_{n \to +\infty} u_n = 0$ donc $u_n^2 = o(u_n)$ et $v_{n+1} - v_n = \frac{1}{u_n} \left( \frac{1}{1-u_n} - 1 \right)$.
On développe : $\frac{1}{1-u} = 1 + u + u^2 + \mathcal{O}(u^3)$
d'où $v_{n+1} - v_n = 1 + u_n + \mathcal{O}(u_n^2) \xrightarrow[n \to +\infty]{} 1$.

D'après Cesaro : $\frac{\sum_{m=0}^{N-1} v_{m+1} - v_m}{N} \xrightarrow[N \to +\infty]{} 1$ d'où $\lim_{n \to +\infty} \frac{v_N}{N} = 1$ et $v_n \sim n$.
Ainsi $u_n \sim \frac{1}{n}$ et $\mathcal{O}(u_n^2) = \mathcal{O}(\frac{1}{n^2})$ tg d'une série abs \boxed{CV}.
Il vient $v_{n+1} - v_n - 1 \sim u_n \sim \frac{1}{n} > 0$ à partir d'un certain rang.
D'après le théorème de sommation des relations de comparaison, comme $\sum \frac{1}{k}$ \boxed{DV} :
\[
\sum_{k=1}^n \underbrace{(v_{k+1} - v_k - 1)}_{v_{n+1} - v_1 - n} \sim \sum_{k=1}^n \frac{1}{k} \sim \ln(n)
\]
Ainsi $v_{n+1} = n + \ln n + o(\ln n)$ et $v_n = n + \ln n + o(\ln n)$ (car $n - (n-1) = o(\ln n)$ et $\ln n - \ln(n-1) = o(\ln n)$).
Il vient $u_n = \frac{1}{n + \ln n + o(\ln n)}$
\[
= \frac{1}{n} \times \frac{1}{1 + \frac{\ln n}{n} + o\left(\frac{\ln n}{n}\right)} = \frac{1}{n} \left( 1 - \frac{\ln n}{n} + o\left(\frac{\ln n}{n}\right) \right) = \frac{1}{n} - \underbrace{\frac{\ln n}{n^2}}_{= \mathcal{O}\left(\frac{1}{n^{3/2}}\right)} + o\left(\frac{\ln n}{n^2}\right)
\]
Ainsi $u_n = \frac{1}{n} + \mathcal{O}\left(\frac{1}{n^{3/2}}\right)$ tg d'une série abs \boxed{CV}.
En reportant : $v_{n+1} - v_n = 1 + \frac{1}{n} + \mathcal{O}\left(\frac{1}{n^{3/2}}\right)$. Il vient alors : $v_{n+1} - v_1 = \sum_{k=1}^n (v_{k+1} - v_k) = n + \underbrace{H_n}_{\sim \ln n + \gamma} + \underbrace{\sum_{k=1}^n \mathcal{O}\left(\frac{1}{k^{3/2}}\right)}_{\to S}$
$v_{n+1} - v_1 = n + \ln n + \gamma + S + o(1)$
Finalement : $\boxed{v_n = n + \ln n + \mathcal{O}(1)}$ [car $n - (n-1) = \mathcal{O}(1)$ et $\ln n - \ln(n-1) = \mathcal{O}(1)$].

\partanswer{c.}
$\boxed{u_n = \frac{1}{n} - \frac{\ln n}{n^2} + o\left(\frac{\ln n}{n^2}\right)}$

\partanswer{d.}
Si $u_0 \in \{0, 1\} : \forall n \geq 1, u_n = 0$.
Si $u_0 < 0$ : Par récurrence : $\forall n \in \N, u_n < 0$, de plus $u_{n+1} < u_n$ : décroissante, strictement négative.
Si $\lim_{n \to +\infty} u_n = l \in \R$ alors $l = 0$, ce qui n'est pas donc $\boxed{\lim_{n \to +\infty} u_n = -\infty}$.
Si $u_0 > 1$, $u_1 < 0$, on en revient au cas précédent.

Rq : Dans ce cas, si on pose $\forall n \in \N, W_n = -u_n$ ; on a $\forall n \in \N, W_{n+1} = -u_{n+1} = -u_n + u_n^2 = W_n + W_n^2 \sim W_n^2$.
On peut espérer que $\exists C > 0 / W_n \sim C^{2^n}$.
Formons $\forall n \in \N, x_n = \frac{\ln W_n}{2^n}$, on a $x_{n+1} = \frac{\ln W_{n+1}}{2^{n+1}} = \frac{\ln(W_n^2 + W_n)}{2^{n+1}} = \frac{\ln W_n^2}{2^{n+1}} + \frac{\ln\left(1 + \frac{1}{W_n}\right)}{2^{n+1}} = \underbrace{\frac{2\ln W_n}{2^{n+1}}}_{x_n} + \underbrace{\frac{\ln\left(1 + \frac{1}{W_n}\right)}{2^{n+1}}}_{\sim \frac{1}{2^{n+1} W_n} = o\left(\frac{1}{2^{n+1}}\right)}$
Ainsi $x_{n+1} - x_n \sim \frac{1}{2^{n+1} W_n} = o\left(\frac{1}{2^{n+1}}\right)$ car $\lim_{n \to +\infty} W_n = +\infty$.
$\sum (x_{n+1} - x_n)$ \boxed{CV} et $(x_n)$ \boxed{CV} vers une limite $K$.
Plus précisément : $W_n = \exp(2^n x_n)$ et $\sum_{k=n}^{+\infty} (x_{k+1} - x_k) = K - x_n = o\underbrace{\left(\sum_{k=n}^{+\infty} \frac{1}{2^{k+1}}\right)}_{= \frac{1}{2^n}}$ car $x_{k+1} - x_k = o\left(\frac{1}{2^{k+1}}\right)$ à partir d'un certain rang.
Ainsi : $2^n x_n = 2^n K + o(1)$ et $\boxed{W_n \sim (e^K)^{2^n}}$.
\end{solution}

\exercise{30}{}

\begin{statementbox}
Soit $u_0 > 0$ et la suite définie par la relation de récurrence $u_{n+1} = u_n + u_n^2$. Mq $\exists C > 0 / u_n \sim C^{2^n}$.
(on pourra former $v_n = \frac{\ln(u_n)}{2^n}$)
\end{statementbox}

\begin{solution}
$f(x) = x + x^2$. $\forall n \in \N, u_{n+1} > u_n$. Par récurrence comme $u_0 > 0$, $\forall n \in \N, u_n > 0$.
Donc la suite admet une limite $l \in \R_+ \cup \{+\infty\}$.
Si $l \in \R$, par continuité de $f : l = l + l^2$ donc $l^2 = 0$ ce qui n'est pas.
Ainsi $\lim_{n \to +\infty} u_n = +\infty$. Et $u_{n+1} = u_n + u_n^2 \sim u_n^2$ $\left( = u_n^2(1 + \frac{1}{u_n}) \right)$.
On pose $v_n = \frac{\ln(u_n)}{2^n}$.
Il vient
\begin{align*}
v_{n+1} &= \frac{\ln(u_{n+1})}{2^{n+1}} = \frac{\ln(u_n^2) + \ln(1+\frac{1}{u_n})}{2^{n+1}} \\
&= \underbrace{\frac{2\ln(u_n)}{2^{n+1}}}_{v_n} + \underbrace{\frac{\ln\left(1+\frac{1}{u_n}\right)}{2^{n+1}}}_{\sim \frac{1}{u_n 2^{n+1}} = o\left(\frac{1}{2^n}\right) \text{ car } \lim_{n \to \infty} u_n = +\infty}
\end{align*}
Ainsi il vient $v_{n+1} - v_n = o\left(\frac{1}{2^{n+1}}\right)$ et $v_{n+1} - v_n > 0$ au delà d'un certain rang. Ainsi $v_{n+1} - v_n \to 0$ donc $(v_{n+1} - v_n)$ \boxed{CV}.
D'après le thm de sommation des relations par comparaison, $\sum_{k=n}^{+\infty} (v_{k+1}-v_k) = o\left( \sum_{k=n}^{+\infty} \frac{1}{2^{k+1}} \right)$ et $\sum (v_{n+1}-v_n)$ \boxed{CV} donc $(v_n)$ converge vers une limite $L \in \R$.
$L - v_n = o\left(\frac{1}{2^n}\right)$ et $v_n = L + o\left(\frac{1}{2^n}\right)$.
Or $u_n = e^{2^n v_n} = (e^L)^{2^n} \times \underbrace{e^{o(1)}}_{\to 1} \sim (e^L)^{2^n}$ d'où le résultat avec $\boxed{C = e^L}$.

\begin{remarkbox}
NB : il ne suffit pas que $\left(\frac{\ln(u_n)}{2^n}\right)_{n \in \N}$ cv vers $L \in \R$ pour que $u_n \sim e^{L 2^n}$. \\
Ex : $u_n = C^{2^n} \times n$.
\end{remarkbox}
\end{solution}

\exercise{48}{}

\begin{statementbox}
Nature de la série $\sum u_n$ avec $\forall n \geq 1$ :
\begin{enumerate}[label=\alph*)]
  \item $u_n = \frac{1}{n^{1+\frac{1}{n}}}$
  \item Soit $(a,b) \in \R^2, u_n = \sin\left(\pi \sqrt{n^2+an+b}\right)$ \quad [$\exists N_0 \in \N, \forall n \geq N_0, n^2+an+b > 0$]
  \setcounter{enumi}{3}
  \item $u_n = e^{- \sum_{k=1}^n \frac{1}{k}}$
  \setcounter{enumi}{5}
  \item $n \geq 2, \quad u_n = (\ln(n \ln n))^{-(\ln n)^\alpha}$
  \item $u_n = \sin\left(\frac{2\pi n!}{e}\right)$ \quad [Rappel : $\frac{1}{e} = \sum_{k=0}^{+\infty} \frac{(-1)^k}{k!}$]
  \item $u_n = \sin(n! \pi e)$
\end{enumerate}
\end{statementbox}

\begin{solution}
\partanswer{a)}
$u_n = \frac{1}{n} \left( \underbrace{\exp\left(-\frac{\ln n}{n}\right)}_{\to 1} \right) \sim \frac{1}{n}$ : $\sum u_n$ \boxed{DV}.

\partanswer{b)}
Pour $n \geq N_0$
$u_n = \sin\left( \pi n \sqrt{1 + \frac{a}{n} + \frac{b}{n^2}} \right), \quad \varepsilon_n = \mathcal{O}\left(\frac{1}{n}\right) \to 0$
On a $\sqrt{1+\varepsilon_n} = 1 + \frac{1}{2}\varepsilon_n - \frac{1}{8}\varepsilon_n^2 + \mathcal{O}(\varepsilon_n^3)$
\begin{align*}
&= 1 + \frac{a}{2n} + \frac{b}{2n^2} - \frac{1}{8}\left( \frac{a^2}{n^2} + \mathcal{O}\left(\frac{1}{n^3}\right) \right) + \mathcal{O}\left(\frac{1}{n^3}\right) \\
&= 1 + \frac{a}{2n} + \frac{1}{n^2} \left( \frac{-a^2 + 4b}{8} \right) + \mathcal{O}\left(\frac{1}{n^3}\right)
\end{align*}
D'où $u_n = \sin\left( \pi n \left( 1 + \frac{a}{2n} + \frac{4b-a^2}{8n^2} + \mathcal{O}\left(\frac{1}{n^3}\right) \right) \right)$
$= \sin\left( n\pi + \frac{\pi a}{2} + \frac{\pi(4b-a^2)}{8n} + \mathcal{O}\left(\frac{1}{n^2}\right) \right)$
$= (-1)^n \sin\left( \frac{\pi a}{2} + \frac{\pi(4b-a^2)}{8n} + \mathcal{O}\left(\frac{1}{n^2}\right) \right)$
Si $a \notin 2\Z : \sin\left( \frac{\pi a}{2} + \frac{\pi(4b-a^2)}{8n} + \mathcal{O}\left(\frac{1}{n^2}\right) \right) \to \sin\frac{\pi a}{2} \neq 0$ et $\sum u_n$ \boxed{DV}.

S'il existe $k \in \Z / a = 2k$.
\begin{align*}
u_n &= (-1)^n (-1)^k \sin\left( \frac{\pi(4b-a^2)}{8n} + \mathcal{O}\left(\frac{1}{n^2}\right) \right) \\
&= (-1)^k \times \underbrace{(-1)^n \frac{\pi(4b-a^2)}{8n}}_{\text{tg d'une série alternée \boxed{CV}}} + \underbrace{\mathcal{O}\left(\frac{1}{n^2}\right)}_{\text{tg d'une série abs \boxed{CV}}} \quad \text{donc } \sum u_n \text{ \boxed{CV}}
\end{align*}

\partanswer{d)}
$u_n = e^{- \sum_{k=1}^n \frac{1}{k}} \sim e^{-\ln(n)} = \frac{1}{n} \xrightarrow[n \to +\infty]{} 0$ ; $\sum \text{DV}$.\\
Par sommation de relations de comparaison : $u_n \sim \frac{1}{n}$, $\sum u_n$ \boxed{\text{DV}}.

\partanswer{f)}
$n \geq 2 \quad u_n = (\ln(n \ln n))^{-(\ln n)^\alpha}$

$u_n = \exp\left( -(\ln n)^\alpha \ln(\ln(n \ln n)) \right)$

$\ln(\ln(n \ln n)) = \ln(\ln n + \ln \ln n) = \ln \ln n + \ln\left(1 + \frac{\ln \ln n}{\ln n}\right) \sim \ln \ln n$

Donc $\exists N_0 \in \N, \forall n \geq N_0, \quad \frac{1}{2} \ln \ln n \leq \ln(\ln n + \ln \ln n) \leq \frac{3}{2} \ln \ln n$

$\forall n \geq N_0, \quad \underbrace{\exp\left(-\frac{3}{2}(\ln n)^\alpha \ln \ln n\right)}_{v_n} \leq u_n \leq \underbrace{\exp\left(-\frac{1}{2}(\ln n)^\alpha \ln \ln n\right)}_{w_n}$

\begin{itemize}
  \item Si $\alpha < 0$ : $\lim_{n \to +\infty} (\ln n)^\alpha \ln \ln n = 0$.
  Donc $\lim v_n = \lim w_n = 1$ et $\lim u_n = 1$. Donc DV grossière.
  
  \item Si $\alpha = 0$ : $\exp\left(-\frac{1}{2}\ln \ln n\right) = \frac{1}{(\ln n)^{1/2}}$ et $\lim_{n \to +\infty} \frac{1}{(\ln n)^{1/2}} = 0$.
  Donc $\sum w_n$ DV donc $\sum u_n$ DV.
  
  \item Si $\alpha > 1$ : $\exists N_1 \in \N, \forall n \geq N_1, \quad (\ln n)^{\alpha-1} \geq 4$\\
  D'où $\forall n \geq \max(N_1, N_0), w_n \leq \exp\left(-2 \ln n\right) = \frac{1}{n^2}$.\\
  Donc $\sum w_n$ CV donc $\sum u_n$ CV.
  
  \item Si $\alpha \in \left]0, 1\right[$ : on a $\frac{1-\alpha}{2} > 0$.\\
  Or $\ln \ln n = o\left( (\ln n)^{\frac{1-\alpha}{2}} \right)$. Donc $\exists N_2 \in \N / \forall n \geq N_2$, 
  $\ln \ln n \leq \frac{2}{3} (\ln n)^{\frac{1-\alpha}{2}}$.\\
  D'où $\forall n \geq N_2, \quad -\frac{3}{2}(\ln n)^\alpha \ln \ln n \geq -(\ln n)^\alpha (\ln n)^{\frac{1-\alpha}{2}} = -(\ln n)^{\frac{1+\alpha}{2}}$ car $\frac{1+\alpha}{2} < 1$.\\
  Alors $\forall n \geq N_2, \quad v_n \geq \exp\left(-(\ln n)^{\frac{1+\alpha}{2}}\right) \geq \frac{1}{n}$.\\
  Donc $\sum v_n$ DV donc $\sum u_n$ DV.
\end{itemize}
Finalement $\sum u_n$ CV ssi $\alpha > 1$.

\textbf{Autre méthode :} on peut essayer d'appliquer un D.A. :\\
On écrit $\ln(\ln(n \ln n)) = \ln(\ln n) + \ln\left( 1 + \frac{\ln \ln n}{\ln n} \right)$
$= \ln \ln n + \frac{\ln \ln n}{\ln n} + \mathcal{O}\left(\left(\frac{\ln \ln n}{\ln n}\right)^2\right)$

$\ln(u_n) = -(\ln n)^\alpha \left( \ln \ln n + \frac{\ln \ln n}{\ln n} + \mathcal{O}\left(\left(\frac{\ln \ln n}{\ln n}\right)^2\right) \right)$
$= -(\ln n)^\alpha \ln \ln n - (\ln n)^{\alpha-1} \ln \ln n + \mathcal{O}\left( \frac{(\ln n)^{\alpha-2}}{(\ln \ln n)^{-2}} \right)$

\partanswer{g)}
$\frac{2\pi n!}{e} = \underbrace{2\pi \sum_{k=0}^n \frac{(-1)^k n!}{k!}}_{\in 2\Z} + \frac{(-1)^{n+1} 2\pi}{n+1} + 2\pi n! \sum_{k=n+2}^{+\infty} \frac{(-1)^k}{k!}$.
Par le CSSA, $\left| \sum_{k=n+2}^{+\infty} \frac{(-1)^k}{k!} \right| \leq \frac{1}{(n+2)!}$, et $2\pi n! \sum_{k=n+2}^{+\infty} \frac{(-1)^k}{k!} = \mathcal{O}\left(\frac{1}{n^2}\right)$
Ainsi $u_n = \sin\left( \frac{(-1)^{n+1} 2\pi}{n+1} + \mathcal{O}\left(\frac{1}{n^2}\right) \right)$
$= \underbrace{(-1)^{n+1} \frac{2\pi}{n+1}}_{\text{tg série alt, \boxed{CV} par CSSA}} + \underbrace{\mathcal{O}\left(\frac{1}{n^2}\right)}_{\text{tg d'une série abs \boxed{CV}}}$.

Et pour $v_n = \sin\left(\frac{\pi n!}{e}\right)$ ??
$\frac{\pi n!}{e} = \pi \underbrace{\sum_{k=0}^{n-2} \frac{(-1)^k n!}{k!}}_{\in 2\Z} + \pi(-1)^{n-1} n + \pi(-1)^n + \frac{(-1)^{n+1}\pi}{n+1} + \mathcal{O}\left(\frac{1}{n^2}\right)$
\begin{align*}
v_n &= (-1)^{n-1} \sin\left( \pi(n-1) + \frac{\pi}{n+1} + \mathcal{O}\left(\frac{1}{n^2}\right) \right) \\
&= + \sin\left( \frac{\pi}{n+1} + \mathcal{O}\left(\frac{1}{n^2}\right) \right) \\
&= \underbrace{\frac{\pi}{n+1}}_{\text{tg d'une série \boxed{DV}}} + \underbrace{\mathcal{O}\left(\frac{1}{n^2}\right)}_{\text{tg d'une série abs \boxed{CV}}}, \quad \text{donc } \sum v_n \text{ \boxed{DV}}
\end{align*}

\partanswer{h)}
$u_n = \sin(n! \pi e) \quad e = \sum_{k=0}^{+\infty} \frac{1}{k!}$

$n! \pi e = \pi \sum_{k=0}^{n-2} \frac{n!}{k!} + \pi n + \pi + \pi n! \sum_{k=n+1}^{+\infty} \frac{1}{k!}$
$= M \pi + \pi n + \pi + \underbrace{\frac{\pi}{n+1} + \mathcal{O}\left(\frac{1}{n^2}\right)}_{= \pi n! \sum_{k=n+1}^{+\infty} \frac{1}{k!}}$ \quad (avec $M \in 2\Z$)

D'où $u_n = \sin\left( \pi(n+1) + \frac{\pi}{n+1} + \mathcal{O}\left(\frac{1}{n^2}\right) \right)$
$= (-1)^{n+1} \sin\left( \frac{\pi}{n+1} + \mathcal{O}\left(\frac{1}{n^2}\right) \right)$
$= \underbrace{(-1)^{n+1} \frac{\pi}{n+1}}_{\text{tg d'une série alternée CV}} + \underbrace{\mathcal{O}\left(\frac{1}{n^2}\right)}_{\text{tg d'une série abs CV}}$

Donc $\sum u_n$ \boxed{\text{CV}}.
\end{solution}

\exercise{49}{}

\begin{statementbox}
\begin{enumerate}
  \item[a.] Montrer la CV et calculer la somme $S = \sum_{m=1}^{+\infty} \left( \sum_{h=m}^{+\infty} \frac{(-1)^{h-1}}{h} \right)$
  \item[j.] Montrer la CV et calculer la somme $\sum_{n=1}^\infty \frac{1}{\binom{n+q}{q+1}}$ où $q \in \N^*$
  \item[2.] Montrer la convergence et calculer la somme de $\sum_{m=1}^{+\infty} \frac{\lfloor m^{1/3} \rfloor - \lfloor (m-1)^{1/3} \rfloor}{4m - m^{1/3}}$
\end{enumerate}
\end{statementbox}

\begin{solution}
\partanswer{a.}
Soit $m \in \N^*, N > m$ :
\begin{align*}
\sum_{h=m}^N \frac{(-1)^{h-1}}{h} &= \sum_{h=m-1}^{N-1} \frac{(-1)^h}{h+1} = \int_0^1 \sum_{h=m-1}^{N-1} (-t)^h dt \\
&= \int_0^1 \frac{(-t)^{m-1} - (-t)^N}{1+t} dt = \int_0^1 \frac{(-t)^{m-1}}{1+t} dt - \underbrace{\int_0^1 \frac{(-t)^N}{1+t} dt}_{|\cdot| \leq \int_0^1 t^N dt = \frac{1}{N+1} \to 0}
\end{align*}
d'où $\sum_{h=m}^{+\infty} \frac{(-1)^{h-1}}{h} = \int_0^1 \frac{(-t)^{m-1}}{1+t} dt$.
Et, soit $N > 1$,
\begin{align*}
\sum_{m=1}^N \left( \sum_{h=m}^{+\infty} \frac{(-1)^{h-1}}{h} \right) &= \int_0^1 \sum_{m=1}^N \frac{(-t)^{m-1}}{1+t} dt \\
&= \int_0^1 \frac{1 - (-t)^N}{(1+t)^2} dt = \int_0^1 \frac{dt}{(1+t)^2} - \underbrace{\int_0^1 \frac{(-t)^N}{(1+t)^2} dt}_{|\cdot| \leq \int_0^1 t^N dt = \frac{1}{N+1} \to 0}
\end{align*}
D'où la \boxed{CV}, et $\sum_{m=1}^{+\infty} \left( \sum_{h=m}^{+\infty} \frac{(-1)^{h-1}}{h} \right) = \int_0^1 \frac{dt}{(1+t)^2} = \left[ -\frac{1}{1+t} \right]_0^1 = \boxed{\frac{1}{2}}$.

\partanswer{j.}
$\forall n \in \N^*, \binom{n+q}{q+1} = \frac{(n+q)!}{(q+1)!(n-1)!} = \frac{(n+q) \times \dots \times n}{(q+1)!} \sim \frac{n^{q+1}}{(q+1)!}$
Ainsi $\frac{1}{\binom{n+q}{q+1}} = \mathcal{O}\left(\frac{1}{n^{q+1}}\right)$ tg d'une série \boxed{CV} d'après Riemann.
Décomposons $\frac{1}{X(X+1)\dots(X+q)} = \sum_{h=0}^q \frac{a_h}{X+h}$
Soit $h \in \llbracket 0, q \rrbracket$. On multiplie par $(X+h)$ et on évalue en $-h$.
$a_h = \frac{1}{(-h)(-h+1)\dots(-h+h-1)(-h+h+1)\dots(-h+q)} = \frac{1}{(-1)^h h! (q-h)!}$
Ainsi $\frac{1}{\binom{n+q}{q+1}} = (q+1)! \sum_{h=0}^q \frac{(-1)^h}{h!(q-h)!} \times \frac{1}{n+h} = (q+1) \sum_{h=0}^q (-1)^h \binom{q}{h} \times \frac{1}{n+h}$

Soit $N > 1$.
\begin{align*}
\sum_{m=1}^N \frac{1}{\binom{m+q}{q+1}} &= (q+1) \sum_{h=0}^q (-1)^h \binom{q}{h} \underbrace{\sum_{m=1}^N \frac{1}{m+h}}_{= H_{h+N} - H_h} \\
&= (q+1) \left[ \sum_{h=0}^q (-1)^h \binom{q}{h} H_{h+N} - \sum_{h=0}^q (-1)^h \binom{q}{h} H_h \right]
\end{align*}
Or pour $h \in \llbracket 0, q \rrbracket$ :
$H_{h+N} = \ln(h+N) + \gamma + o(1) = \ln N + \gamma + \ln\left(1 + \frac{h}{N}\right) + o(1) = \ln N + \gamma + o(1)$
D'où $\sum_{m=1}^N \frac{1}{\binom{m+q}{q+1}} = (q+1) \left[ (\ln N + \gamma) \underbrace{\sum_{h=0}^q (-1)^h \binom{q}{h}}_{(1-1)^q = 0 \text{ car } q \in \N^*} + o(1) - \sum_{h=0}^q (-1)^h \binom{q}{h} H_h \right]$
\[
\xrightarrow[N \to +\infty]{} (q+1) \sum_{h=0}^q (-1)^{h+1} \binom{q}{h} H_h
\]

\partanswer{2.}
Soit $(k,m) \in \N^2$, on a $k = \lfloor m^{1/3} \rfloor$ ssi $k \leq m^{1/3} < k+1$ ssi $k^3 \leq m < (k+1)^3$.
Si $\forall k \in \N, m \neq k^3$ : $k^3 < m < (k+1)^3$
donc $k^3 \leq m-1 < (k+1)^3$
et $\lfloor (m-1)^{1/3} \rfloor = \lfloor m^{1/3} \rfloor$.

Or si $m = k^3$ : $\lfloor m^{1/3} \rfloor = k$, $\lfloor (m-1)^{1/3} \rfloor = k-1$.
Ainsi $\lfloor m^{1/3} \rfloor - \lfloor (m-1)^{1/3} \rfloor = \begin{cases} 1 & \text{si } m = k^3 \\ 0 & \text{sinon} \end{cases}$.

Il vient (dans $\R_+ \cup \{+\infty\}$) : $\sum_{m=1}^{+\infty} u_m = \sum_{k=1}^{+\infty} \frac{1}{4k^3 - k}$ CV 
car $\frac{1}{4k^3 - k} = \mathcal{O}\left(\frac{1}{k^3}\right)$.

On décompose en éléments simples $\frac{1}{4X^3 - X} = \frac{1}{X(4X^2 - 1)} = \frac{1}{X(2X-1)(2X+1)}$
$= \frac{a}{X} + \frac{b}{2X-1} + \frac{c}{2X+1}$
On a $\begin{cases} a = -1 \\ b = 1 \\ c = 1 \end{cases}$

Soit $N \geq 1$ : $\sum_{k=1}^N \frac{1}{4k^3 - k} = \sum_{k=1}^N -\frac{1}{k} + \sum_{k=1}^N \frac{1}{2k-1} + \sum_{k=1}^N \frac{1}{2k+1}$
\[
= -H_N + \underbrace{1 + \frac{1}{3} + \dots + \frac{1}{2N-1}}_{H_{2N} - \frac{1}{2}H_N} + \underbrace{\frac{1}{3} + \frac{1}{5} + \dots + \frac{1}{2N+1}}_{H_{2N} - 1 - \frac{1}{2}H_N + \frac{1}{2N+1}}
\]
\end{solution}
